% ============================================================
%  REVIEW COMMENTS & REVISION SUGGESTIONS
%  Paper: "Dropout-Aware Training for Sensor-Robust Fault
%          Localization with Spatio-Temporal GNNs"
%  Reviewer: Dr. M Anbazhagan (Supervisor)
%  Prepared for: P K Sangameswar (M.Tech. Student)
%  Target venue: Procedia Computer Science (ICMLDE 2025)
%  Page budget: <= 12 pages
% ============================================================

\documentclass[11pt,a4paper]{article}

\usepackage[margin=2.5cm]{geometry}
\usepackage{xcolor}
\usepackage{mdframed}
\usepackage{enumitem}
\usepackage{booktabs}
\usepackage{hyperref}
\usepackage{tcolorbox}
\usepackage{titlesec}
\usepackage{parskip}

% ---- Colour palette for comment boxes ----------------------
\definecolor{critred}{RGB}{192,0,0}
\definecolor{warnorange}{RGB}{204,102,0}
\definecolor{suggestblue}{RGB}{0,70,127}
\definecolor{okgreen}{RGB}{0,112,60}
\definecolor{trimviolet}{RGB}{102,0,153}
\definecolor{lightgray}{RGB}{245,245,245}

% ---- Comment box environments ------------------------------
\tcbuselibrary{skins,breakable}

\newtcolorbox{critbox}[1]{
  breakable, colback=red!5, colframe=critred,
  fonttitle=\bfseries\color{critred},
  title={[CRITICAL] #1}}

\newtcolorbox{warnbox}[1]{
  breakable, colback=orange!8, colframe=warnorange,
  fonttitle=\bfseries\color{warnorange},
  title={[REVISION NEEDED] #1}}

\newtcolorbox{suggestbox}[1]{
  breakable, colback=blue!5, colframe=suggestblue,
  fonttitle=\bfseries\color{suggestblue},
  title={[SUGGESTION] #1}}

\newtcolorbox{okbox}[1]{
  breakable, colback=green!5, colframe=okgreen,
  fonttitle=\bfseries\color{okgreen},
  title={[STRENGTH] #1}}

\newtcolorbox{trimbox}[1]{
  breakable, colback=violet!5, colframe=trimviolet,
  fonttitle=\bfseries\color{trimviolet},
  title={[TRIM TO FIT 12pp] #1}}

% ---- Inline annotation commands ----------------------------
\newcommand{\addref}{\textcolor{critred}{\textbf{[CITE?]}}}
\newcommand{\clarify}[1]{\textcolor{warnorange}{\textit{[Clarify: #1]}}}
\newcommand{\reword}[1]{\textcolor{suggestblue}{\textit{[Reword: #1]}}}
\newcommand{\trimhere}{\textcolor{trimviolet}{\textbf{$\blacktriangleleft$\,TRIM}}}

% ---- Section title style -----------------------------------
\titleformat{\section}{\large\bfseries\color{suggestblue}}{}{0em}{}[\titlerule]
\titleformat{\subsection}{\normalsize\bfseries}{}{0em}{}

% ============================================================
\begin{document}

\begin{center}
  {\Large\bfseries Review Notes \& Revision Guide}\\[4pt]
  {\large\itshape Dropout-Aware Training for Sensor-Robust Fault Localization\\
  with Spatio-Temporal Graph Neural Networks}\\[6pt]
  % {\normalsize Procedia Computer Science \textemdash{} ICMLDE 2025 \quad
  % Page budget: \textbf{$\leq$ 12 pages}}\\[2pt]
  % {\small Reviewed by: Dr.\ M.\ Anbazhagan \quad Date: August 2026}
\end{center}

\vspace{6pt}
\hrule
\vspace{10pt}

% ============================================================
\section{Overall Assessment}
% ============================================================

\begin{okbox}{Core Contributions Are Solid}
The paper makes three clearly differentiated, internally consistent contributions:
(i)~a first systematic sensor-dropout robustness study for STGNN-based fault
localisation; (ii)~dropout-aware training as a simple, architecture-agnostic
mitigation; and (iii)~the SDR-STGNN ablation that cleanly isolates training-time
vs.\ architectural effects. The multi-seed, shared-mask evaluation protocol is a
particular strength, as it ensures paired comparisons. The findings (AUDC rising
from 0.414 to 0.916; retention at $k{=}10$ from 7\% to 81\%) are both
statistically meaningful and practically significant.
\end{okbox}

\begin{warnbox}{Page Count Must Be Reduced}
The draft as received renders at approximately \textbf{15--16 pages} in the
Procedia two-column template.  The ICMLDE call-for-papers caps regular papers at
\textbf{12 pages} (including references and figures). Detailed trimming
instructions are given section by section below in
\textcolor{trimviolet}{\textbf{[TRIM]} boxes}. The estimated saving from all
recommended cuts is \textbf{3.5--4 pages}, which will bring the paper within
budget without sacrificing any primary result.
\end{warnbox}

% ============================================================
\section{Abstract}
% ============================================================

\begin{okbox}{Well-structured; leads with the practical problem}
The abstract correctly motivates partial observability, positions the gap
(no prior STGNN robustness study), and quantifies the headline result. The
progression from problem $\to$ gap $\to$ method $\to$ result $\to$ conclusion is
clean.
\end{okbox}

\begin{suggestbox}{Minor wording tightening}
\begin{enumerate}[leftmargin=*,label=\arabic*.]
  \item \textbf{Line 3 -- redundant phrase.}
    ``\ldots yet it remains difficult under partial observability, \emph{where
    only a limited number of micro-phasor measurement units ($\mu$PMUs) monitor
    networks of hundreds of buses}.'' The parenthetical definition of
    $\mu$PMUs is repeated verbatim in the Introduction. Keep it in the abstract
    and shorten the Introduction's first mention to a back-reference:
    ``\ldots micro-phasor measurement units ($\mu$PMUs; see Abstract)'',
    or simply rely on the abstract's definition.
  \item \textbf{Quantify the clean-accuracy cost explicitly.}
    Add one number: ``\ldots at virtually no cost to clean performance
    ($\Delta$F1 $<$ 0.01 at $k{=}0$).'' This is a key selling point that
    reviewers may miss if buried in Table~6.
  \item \textbf{Acronym SDR-STGNN.} Do not introduce the acronym in the
    abstract if it is not used again there. Either drop the parenthetical
    ``(SDR-STGNN)'' or use it in the very next sentence.
\end{enumerate}
\end{suggestbox}

% ============================================================
\section{Section 1 -- Introduction}
% ============================================================

\begin{warnbox}{Four contributions listed; paragraph structure can be tightened}
\begin{enumerate}[leftmargin=*,label=\arabic*.]
  \item \textbf{Bullet~1 (dataset + baselines)} is essentially a replication
    study step, not a novel contribution. Re-label it as ``Experimental
    Infrastructure'' or integrate it into the setup description in
    Section~3 to strengthen the contribution list. Reviewers count bullets;
    four dilutes the three genuinely novel ones.
  \item \textbf{Paragraph 2} (``Graph neural networks are a natural
    fit\ldots'') and \textbf{Paragraph 3} (describing~\cite{Karabulut2026})
    together spend $\sim$180 words on background that is better placed in
    Section~2. Move the summary of~\cite{Karabulut2026} to Related Work
    Table~1 row and trim this paragraph to two sentences.
  \item \textbf{Final roadmap paragraph} (``The remainder of this
    paper\ldots'') is a standard venue boilerplate and consumes $\sim$60
    words. For a 12-page paper at this venue, one sentence suffices or it
    may be omitted entirely.
\end{enumerate}
\end{warnbox}

\begin{trimbox}{Introduction: estimated saving $\approx$ 0.5 page}
\begin{itemize}[leftmargin=*]
  \item Remove the five-sentence description of \cite{Karabulut2026}'s
    architecture (``\ldots GRU for learning temporal patterns\ldots 94.7--94.9'')
    from the Introduction; it duplicates Table~1 and Section~3.2.
    Replace with a single sentence plus a forward reference:
    \textit{``Building on the STGNN framework of~\cite{Karabulut2026} (see
    Section~3.2 and Table~1), we extend \ldots''}
  \item Collapse the roadmap to: \textit{``Section~2 reviews related work;
    Sections~3--4 describe the methodology and setup; Sections~5--6 present
    results; Sections~7--8 discuss and conclude.''}
\end{itemize}
\end{trimbox}

% ============================================================
\section{Section 2 -- Related Work}
% ============================================================

\begin{okbox}{Table 1 is an effective positioning device}
The seven-column table gives reviewers an at-a-glance comparison that clearly
shows the unique combination of \{sparse obs.\ + dropout @ test + mitigation\}
only this work addresses. Keep it.
\end{okbox}

\begin{warnbox}{Prose around Table 1 is thin}
There is currently no explanatory paragraph that \emph{interprets} why existing
methods fall short beyond what the checkmarks convey. Add 2--3 sentences
below the table that draw a conclusion:
\begin{quote}
\textit{``The table reveals that no prior method simultaneously addresses sparse
observability and test-time sensor dropout with active mitigation. Imputation
approaches~\cite{Zhang2025,Fan2024} assume a fixed, known missingness pattern
incompatible with random inference-time failures. This gap motivates the
evaluation and mitigation study presented here.''}
\end{quote}
\end{warnbox}

\begin{suggestbox}{Citation completeness}
\begin{enumerate}[leftmargin=*,label=\arabic*.]
  \item \cite{Karabulut2025} (\texttt{arXiv:2510.03571}) appears in the
    reference list but is not cited in the body text. Either cite it
    explicitly in Section~2 as a companion generalization study, or remove it
    from the reference list to save space.
  \item The paper claims ``to the best of our knowledge, no previous work has
    studied how STGNN-based fault localization is affected when some sensors
    fail during inference.'' This is a strong claim. Recommend a brief
    search of 2024--2025 literature on \textit{missing-sensor robustness in
    power systems} and IEC 61850 communication-failure studies before
    submission to ensure it holds.
\end{enumerate}
\end{suggestbox}

% ============================================================
\section{Section 3 -- Materials and Methods}
% ============================================================

\begin{okbox}{Section 3 is the paper's methodological core and is well-written}
The dataset generation, the measured-only graph formulation, and the dropout-aware
training algorithm (Algorithm~1) are described with sufficient reproducibility
detail. The six-dimensional node feature including the availability mask is a
technically elegant design choice that is clearly motivated.
\end{okbox}

\begin{warnbox}{Section 3.1 -- dataset description is over-detailed for a conference paper}
Two specific issues:
\begin{enumerate}[leftmargin=*,label=\arabic*.]
  \item The paragraph on \textbf{voltage response and Fig.~2} restates the
    obvious (``low-impedance faults produce pronounced voltage drops'').
    Retain the figure but cut the explanatory paragraph to one sentence:
    \textit{``Fig.~2 illustrates that high-impedance (10\,$\Omega$) faults
    produce shallow voltage sags ($<$0.02\,p.u.), making them the hardest
    class to localise.''}
  \item Table~2 is comprehensive but the \textbf{window-level distribution
    sub-table} (train/val/test counts) consumes vertical space disproportionate
    to its informational value. Collapse the three split rows into a single
    inline sentence:
    \textit{``The 1.65\,M windows are partitioned 70/15/15 at the event
    level (train: 1,155\,k; val: 248\,k; test: 247\,k).''}
    and remove the lower half of Table~2.
\end{enumerate}
\end{warnbox}

\begin{warnbox}{Section 3.2 -- five baselines described redundantly}
The descriptions of RGCN, RGSAGE-Mean/Max, and RGATv2 faithfully recapitulate
\cite{Karabulut2026}. For a paper whose contribution is robustness evaluation,
not architecture design, condense to a single paragraph:
\begin{quote}
\textit{``We adopt the five STGNN variants of~\cite{Karabulut2026}: a temporal-only
GRU, RGCN~\cite{Hamilton2017}, RGSAGE-Mean, RGSAGE-Max, and
RGATv2~\cite{Brody2022}. All share a GRU temporal encoder (128 hidden units)
followed by an optional spatial stage over the 25-node measured-only graph; a
shared classification head and soft voting produce the graph-level prediction.
We extend the spatial depth to four hops (Section~5) and otherwise adopt the
architectures unchanged.''}
\end{quote}
This saves approximately 0.4 page without losing reproducibility, since the
operators are standard and cited.
\end{warnbox}

\begin{suggestbox}{Section 3.3 -- robustness metrics: define AUDC with a cleaner integral notation}
The current continuous integral form of AUDC,
\[
  \mathrm{AUDC} = \frac{1}{k_{\max}-k_{\min}}\int_{k_{\min}}^{k_{\max}} R(k)\,dk,
\]
implies a continuous curve, but $k$ is discrete. Add the sentence: \textit{``In
practice we evaluate AUDC via trapezoidal integration over
$k \in \{0,1,2,5,10\}$.''} (already in the text but appears three lines after
the integral -- move it immediately after.)
\end{suggestbox}

\begin{warnbox}{Section 3.5 -- SDR-STGNN: Equations 2 and 3 need labelling consistency}
Equation~(1) defines $a_i$; Equations~(2) and~(3) define gated aggregation.
Currently the text references ``Eq.~3'' and ``Eq.~3.5'' interchangeably.
Standardise to \texttt{\textbackslash eqref} throughout and verify that the
equation counter resets correctly in the Procedia template.
\end{warnbox}

% ============================================================
\section{Section 4 -- Experimental Setup}
% ============================================================

\begin{trimbox}{Section 4 can be reduced by $\approx$ 0.3 page}
\begin{itemize}[leftmargin=*]
  \item The \textbf{hardware paragraph} (NVIDIA T1000, Python 3.12, PyTorch
    2.5.1, CUDA 12.1, PyG 2.8, NumPy 2.5, scikit-learn 1.9,
    OpenDSSDirect 0.9.4) is important for reproducibility but reads like a
    version-number list. Move the exact versions to a footnote or a
    \textit{Reproducibility Statement} at the end (some venues now prefer
    this), and retain only: \textit{``All experiments were run on a single
    NVIDIA T1000 GPU (8\,GB) using PyTorch 2.5.1 with PyTorch Geometric 2.8.''}
  \item The \textbf{``Dropout-aware training'' paragraph} in Section~4
    repeats what was said in Section~3.4. Remove it entirely; a
    cross-reference suffices: \textit{``Dropout-aware training follows
    Section~3.4 with no additional hyperparameter tuning.''}
\end{itemize}
\end{trimbox}

% ============================================================
\section{Section 5 -- Differences from the Base Paper}
% ============================================================

\begin{trimbox}{Section 5 is the single largest trimming opportunity ($\approx$ 1 full page)}
This section exists to pre-empt reviewer questions about why the reported F1
(0.85) differs from \cite{Karabulut2026}'s (0.947). That is a legitimate
concern, but it does not require a dedicated section in a 12-page paper.
Recommended action:
\begin{enumerate}[leftmargin=*,label=\arabic*.]
  \item \textbf{Fold into Section~3 or Section~4.} Add a clearly marked
    paragraph titled ``\textit{Deviations from~\cite{Karabulut2026}}'' at the
    end of Section~3.2 (baselines) and list all four differences in a compact
    two-column table:
    \begin{center}
    \small
    \begin{tabular}{lll}
      \toprule
      Aspect & \cite{Karabulut2026} & This work \\
      \midrule
      Spatial depth & 1 hop & 4 hops \\
      Training loss & Class-weighted CE & Unweighted CE \\
      Corpus size & Not disclosed & 1.65\,M windows (50 runs/scenario)\\
      Green topology & Evaluated & Not evaluated \\
      \bottomrule
    \end{tabular}
    \end{center}
    Follow with a single sentence: \textit{``All comparisons in this work are
    conducted on our own corpus and splits, so robustness effects are
    internally consistent regardless of absolute F1 differences from the base
    paper.''}
  \item \textbf{Remove the prose paragraphs} that justify each difference at
    length (Spatial depth, Training Loss, Corpus size, Scope not covered).
    The table above communicates the same information in $<$10 lines instead
    of $\sim$60.
\end{enumerate}
\textbf{Estimated saving: 0.9--1.0 page.}
\end{trimbox}

% ============================================================
\section{Section 6 -- Results}
% ============================================================

\begin{okbox}{The results section is the paper's strongest part}
The three-stage structure (baseline accuracy $\to$ degradation curves $\to$
mitigation comparison) follows a logical progression. Tables~3, 5, and~6 are
clean, informative, and well-captioned. Figures~3 and~4 are publication-quality
and directly support the key claims.
\end{okbox}

\begin{warnbox}{Section 6.1 -- Table 3 caption and spatial-depth Table 4}
\begin{enumerate}[leftmargin=*,label=\arabic*.]
  \item \textbf{Table~4} (depth sweep) is presented as ``diagnostic'' but is
    given nearly the same visual weight as Table~3. Demote it: either move it
    to an appendix (if the venue allows) or shrink it to an inline sentence:
    \textit{``A single-seed depth sweep on RGSAGE-Max shows macro-F1
    increasing monotonically from 0.509 (1 hop) to 0.846 (4 hops, our
    operating point).''}
  \item \textbf{RGATv2 training instability} ($\pm$0.087) is mentioned in
    the results narrative but only cautioned against rather than
    investigated. Add a one-sentence recommendation: \textit{``Stabilising
    RGATv2 at depth~4 (e.g., via attention dropout tuning or warmup
    scheduling) is left as future work.''} This proactively addresses an
    inevitable reviewer question.
\end{enumerate}
\end{warnbox}

\begin{warnbox}{Section 6.2 -- Narrative repeats Table 5 numbers verbatim}
The paragraph beginning \textit{``RGSAGE-Max, the strongest model under full
observability, is also the least robust\ldots''} re-states every number already
in Table~5. Convert to a higher-level interpretation:
\begin{quote}
\textit{``Table~5 reveals a clear inversion: accuracy rank (RGSAGE-Max $>$
RGSAGE-Mean $>$ RGCN $>$ RGATv2 $>$ GRU) is exactly reversed by robustness
rank (RGATv2 $\approx$ RGCN $>$ RGSAGE-Mean $>$ GRU $>$ RGSAGE-Max). No
baseline retains more than 32\% of its full-observability F1 at 40\% sensor
loss. The mechanism is the aggregation operator: max-pooling allows a single
corrupted (zeroed) neighbour embedding to dominate a live node's update,
while mean and normalised-sum dilute it. This motivates both mitigations
in Section~6.3.''}
\end{quote}
\textbf{Estimated saving: 0.25 page.}
\end{warnbox}

\begin{suggestbox}{Section 6.3 -- Report per-fault-resistance robustness for dropout-aware model}
Table~6 reports only aggregate macro-F1. Since Table~3 shows that
\textbf{10\,$\Omega$ high-impedance faults} are the hardest class (F1 drops
from 0.949 to 0.676 as resistance increases), readers will naturally ask: does
dropout-aware training help \emph{equally} across fault resistances, or does it
fail first on high-impedance faults?\\[4pt]
Add a single per-resistance row to Table~6, or state explicitly:
\textit{``Dropout-aware training preserves per-resistance ranking: at $k{=}10$
the 10\,$\Omega$ class retains $X$\% of its full-sensor F1, compared with
$Y$\% for the 0.1\,$\Omega$ class (see supplementary Table~S1).''} Even a
qualitative statement prevents a reviewer rejection on this point.
\end{suggestbox}

% ============================================================
\section{Section 7 -- Discussion}
% ============================================================

\begin{trimbox}{Section 7.1 is the second largest trimming target ($\approx$ 0.7 page)}
Section~7.1 (\textit{``Which Architectures Are Most Robust to Sensor Loss''})
re-explains, in prose, the results already visible in Tables~5 and~6 and
Figures~3 and~4. The paragraph beginning ``This ordering follows from how each
operator treats the zeroed embedding\ldots'' is insightful \emph{mechanistic
reasoning} and must be kept. Everything else that re-narrates numbers should be
cut.

\medskip
\textbf{Recommended structure for a trimmed Section~7:}
\begin{enumerate}[leftmargin=*,label=\arabic*.]
  \item Two-sentence summary of the accuracy-vs-robustness inversion (keep
    the mechanistic explanation of max-pooling fragility).
  \item One paragraph on why dropout-aware training works (currently split
    across 7.1 and the caption of Fig.~4 -- unify here).
  \item One paragraph on the SDR-STGNN finding (gating removes multi-hop
    rerouting paths, hence slightly lower AUDC than ungated -- this is
    already well-stated in Section~6.3 and need not be repeated).
\end{enumerate}
The Limitations (Section~7.2) are complete and well-scoped -- keep as-is.
\textbf{Estimated saving: 0.5--0.7 page.}
\end{trimbox}

\begin{warnbox}{Missing discussion point: practical deployment implications}
The paper demonstrates that a model trained with dropout-aware augmentation
is robust to 40\% sensor loss. This has a direct practical implication for
$\mu$PMU network operators: \emph{what minimum sensor availability guarantees
acceptable service?} Add a two-sentence design guideline:
\begin{quote}
\textit{``For the IEEE 123-bus feeder with 25 $\mu$PMUs, dropout-aware training
maintains macro-F1 $>$ 0.67 even at 40\% loss ($k{=}10$). This suggests that
operators can tolerate losing up to $\sim$10 of 25 sensors during scheduled
maintenance without degrading the localisation system below an operationally
useful threshold, provided the model is retrained with dropout-aware
augmentation.''}
\end{quote}
This connects the academic result to SDG~7 (Affordable and Clean Energy) and
SDG~11 (Sustainable Cities), and strengthens the paper's impact statement.
\end{warnbox}

% ============================================================
\section{Section 8 -- Conclusion and Future Work}
% ============================================================

\begin{okbox}{Conclusion is well-scoped and appropriately brief}
The conclusion correctly identifies the three contributions, quantifies the
headline result, and proposes four concrete directions for future work (correlated
failures, larger feeders, cross-topology generalisation, adaptive gating). No
changes needed for content.
\end{okbox}

\begin{suggestbox}{Minor: cite the base paper in the opening sentence}
The conclusion opens with ``Building on the STGNN fault-localization methodology
of~[1]'' -- this is correct and should be retained. However, the citation style
should match the body text: use \texttt{\textbackslash cite\{Karabulut2026\}}
rather than the manual ``[1]'' in the final manuscript to ensure the bibliography
auto-numbers correctly in the Procedia template.
\end{suggestbox}

% ============================================================
\section{References}
% ============================================================

\begin{warnbox}{Reference [2] / \texttt{Karabulut2025} is uncited in the body}
Reference~[2] (Karabulut et al., arXiv:2510.03571, ``Generalization of GNN
models for distribution grid fault detection'') appears in the list but is never
cited. Either:
\begin{itemize}[leftmargin=*]
  \item Cite it in Section~2 as the companion generalisation study:
    \textit{``Karabulut et al.~\cite{Karabulut2025} separately study
    cross-feeder generalisation; our work is complementary, addressing
    sensor-dropout robustness on a fixed topology.''}
  \item Or remove it to save one reference slot (Procedia does not enforce a
    minimum, but an uncited reference raises reviewer eyebrows).
\end{itemize}
\end{warnbox}

\begin{suggestbox}{Verify arXiv references before final submission}
References~[1] and~[2] are arXiv preprints. Before camera-ready submission,
confirm whether \cite{Karabulut2026} has appeared in a peer-reviewed venue
(IEEE Transactions on Smart Grid, etc.) and update accordingly. Conference
proceedings rank preprint citations lower.
\end{suggestbox}

% ============================================================
\section{Figures and Tables}
% ============================================================

\begin{okbox}{Figures 3 and 4 are publication-ready}
Both figures use consistent colours, labelled axes, error bars ($\pm$1~std),
and clear legends. The non-linear x-axis ($k \in \{0,1,2,5,10\}$) correctly
emphasises the low-dropout behaviour where differences are most interpretable.
No changes needed.
\end{okbox}

\begin{warnbox}{Figure 1 resolution and caption}
Fig.~1 (the measured-only graph) appears visually adequate but:
\begin{enumerate}[leftmargin=*,label=\arabic*.]
  \item Export at $\geq$300\,dpi as PDF or EPS (not PNG) for the Procedia
    camera-ready to avoid pixelation in print.
  \item The caption states ``Red rings mark the 6 $\mu$PMU buses that are
    also candidate fault locations''. Confirm this matches the actual figure
    (some red rings may be partially obscured by node labels in the
    submitted PDF).
\end{enumerate}
\end{warnbox}

\begin{suggestbox}{Consider adding a compact architecture diagram}
A small schematic (one column width, $\leq$6\,cm tall) showing the GRU $\to$
GNN $\to$ classifier $\to$ soft-vote pipeline, with the dropout-aware masking
step indicated as an annotation on the input, would help readers who come from
a power-systems rather than a GNN background. This is optional but would
strengthen the paper's accessibility.
\end{suggestbox}

% ============================================================
\section{Summary: Page-Saving Roadmap}
% ============================================================

\begin{tcolorbox}[colback=violet!6,colframe=trimviolet,
  fonttitle=\bfseries\color{trimviolet},
  title={Trimming Plan to Reach $\leq 12$ Pages}]

\begin{tabular}{p{0.55\linewidth}r}
\toprule
\textbf{Action} & \textbf{Est.\ saving} \\
\midrule
\S1: Remove base-paper architecture recap + shorten roadmap
  & $\sim$0.5 pp \\
\S3.1: Collapse Table~2 split sub-table to one sentence
  & $\sim$0.15 pp \\
\S3.1: Trim voltage-sag paragraph to one sentence
  & $\sim$0.10 pp \\
\S3.2: Condense 5-baseline descriptions to one paragraph
  & $\sim$0.4 pp \\
\S4: Remove duplicate dropout-aware paragraph; footnote versions
  & $\sim$0.25 pp \\
\S5: Replace with compact deviation table (biggest single saving)
  & $\sim$0.9 pp \\
\S6.1: Demote Table~4 to inline sentence
  & $\sim$0.15 pp \\
\S6.2: Replace number-recitation with interpretive paragraph
  & $\sim$0.25 pp \\
\S7.1: Trim to mechanistic explanation only
  & $\sim$0.55 pp \\
\midrule
\textbf{Total estimated saving} & $\mathbf{\sim 3.25}$\,\textbf{pp} \\
\bottomrule
\end{tabular}

\medskip
\textit{Starting from $\approx$15.5 pages, the cuts above bring the manuscript
to $\approx$12.2 pages. A further 0.2-page saving is easily achieved by reducing
inter-paragraph spacing and ensuring all figures are cropped tightly, which the
Procedia template handles automatically if bounding boxes are set correctly.}
\end{tcolorbox}

% ============================================================
\section{Priority Action List for Student}
% ============================================================

\begin{tcolorbox}[colback=red!5,colframe=critred,
  fonttitle=\bfseries\color{critred},
  title={Must-Do Before Resubmission (in priority order)}]
\begin{enumerate}[leftmargin=*]
  \item \textbf{Trim Section~5} (Differences from Base Paper) to a compact
    table of four rows -- single largest page saving.
  \item \textbf{Add per-resistance robustness data} for the dropout-aware
    model (Table~6 extension or supplementary) -- closes an obvious reviewer gap.
  \item \textbf{Cite or remove} Reference~[2] (\texttt{Karabulut2025}).
  \item \textbf{Condense Section~3.2} baseline descriptions to one paragraph.
  \item \textbf{Add practical deployment implication} paragraph in Section~7
    (connects to SDG~7/11 and strengthens impact).
  \item \textbf{Verify} arXiv references~[1] and~[2] for peer-reviewed
    publication status.
  \item \textbf{Export} Fig.~1 at $\geq$300\,dpi as vector PDF.
  \item \textbf{Fix equation cross-reference style} in Section~3.5
    (use \texttt{\textbackslash eqref} consistently).
\end{enumerate}
\end{tcolorbox}

% ============================================================
% Bibliography stub (matches paper's reference list)
% ============================================================
\begin{thebibliography}{99}

\bibitem{Karabulut2026}
B.~Karabulut, C.~Manna, and C.~Develder,
``Robustness of spatio-temporal graph neural networks for fault location in
partially observable distribution grids,''
\textit{arXiv preprint arXiv:2604.20403}, 2026.

\bibitem{Karabulut2025}
B.~Karabulut, C.~Manna, and C.~Develder,
``Generalization of graph neural network models for distribution grid fault
detection,''
\textit{arXiv preprint arXiv:2510.03571}, 2025.

\bibitem{Hamilton2017}
W.~Hamilton, Z.~Ying, and J.~Leskovec,
``Inductive representation learning on large graphs,''
\textit{Advances in Neural Information Processing Systems}, vol.~30, 2017.

\bibitem{Brody2022}
S.~Brody, U.~Alon, and E.~Yahav,
``How attentive are graph attention networks?''
\textit{ICLR}, 2022.

\bibitem{Zhang2025}
L.~Zhang et al.,
``Fault location in power distribution networks with massive missing data,''
\textit{IEEE Trans.\ Autom.\ Sci.\ Eng.}, vol.~22, pp.~16825--16837, 2025.

\bibitem{Fan2024}
M.~Fan et al.,
``Fault location method of distribution network based on VGAE-GraphSAGE,''
\textit{Processes}, vol.~12, no.~10, p.~2179, 2024.

\end{thebibliography}

\end{document}